% !TEX root = ../../ctfp-print.tex

\lettrine[lhang=0.17]{自}{由構成は随伴の強力な応用だ。}\newterm{自由関手}（free functor）は\newterm{忘却関手}（forgetful functor）への左随伴として定義される。忘却関手は通常、何らかの構造を「忘れる」非常に単純な関手だ。たとえば、多くの興味深い圏は集合の上に構築されている。しかし、それらの集合を抽象化する圏論的対象には内部構造がない――要素を持たない。それでも、これらの対象はしばしば集合の記憶を持っており、与えられた圏$\cat{C}$から$\Set$への写像――関手――が存在するという意味においてそうだ。$\cat{C}$のある対象に対応する集合は、その\newterm{台集合}（underlying set）と呼ばれる。

モノイドはそのような台集合――要素の集合――を持つ対象だ。モノイドの圏$\cat{Mon}$から集合の圏への忘却関手$U$があり、モノイドをその台集合に写す。また、モノイド射（準同型）を集合間の関数に写す。

私は$\cat{Mon}$を分裂した人格を持つと考えたい。一方では、乗算と単位元を持つ集合の集まりだ。他方では、唯一の構造が射に符号化された特徴のない対象を持つ圏だ。乗算と単位を保存する集合関数はすべて$\cat{Mon}$の射を生む。\\
\newline
覚えておくべきこと：

\begin{itemize}
  \tightlist
  \item
        同じ集合に写されるモノイドが複数ある場合があり、
  \item
        台集合間の関数より少ない（または多くて同じ数の）モノイド射が存在する。
\end{itemize}

\noindent
忘却関手$U$の左随伴である関手$F$は、生成集合から自由モノイドを構築する自由関手だ。随伴は、以前に説明した自由モノイドの普遍的構成から導かれる。\footnote{第13章の\hyperref[free-monoids]{自由モノイド}を参照。}

\begin{figure}[H]
  \centering
  \includegraphics[width=0.6\textwidth]{images/forgetful.jpg}
  \caption{モノイド$m_1$と$m_2$は同じ台集合を持つ。$m_2$と$m_3$の台集合間の関数は、それらの間のモノイド射よりも多く存在する。}
\end{figure}

\noindent
hom集合の観点では、この随伴を次のように書ける：
\[\cat{Mon}(F x, m) \cong \Set(x, U m)\]
この（$x$と$m$について自然な）同型は次のことを示す：

\begin{itemize}
  \tightlist
  \item
        $x$によって生成される自由モノイド$F x$と任意のモノイド$m$の間のすべてのモノイド準同型に対して、生成元の集合$x$を$m$の台集合に埋め込む一意な関数が存在する。それは$\Set(x, U m)$内の関数だ。
  \item
        $x$をある$m$の台集合に埋め込むすべての関数に対して、$x$によって生成された自由モノイドとモノイド$m$の間の一意なモノイド射が存在する。（これは普遍的構成で$h$と呼んだ射だ。）
\end{itemize}

\begin{figure}[H]
  \centering
  \includegraphics[width=0.8\textwidth]{images/freemonadjunction.jpg}
\end{figure}

\noindent
直感的には、$F x$は$x$を基盤として構築できる「最大」のモノイドだ。モノイドの内部を見ることができれば、$\cat{Mon}(F x, m)$に属するすべての射がこの自由モノイドを何らかの他のモノイド$m$に\emph{埋め込む}ことがわかるだろう。それはいくつかの要素を同一視することで行われる。特に、$F x$の生成元（つまり$x$の要素）を$m$に埋め込む。随伴は、右辺の$\Set(x, U m)$からの関数によって与えられる$x$の埋め込みが、左辺のモノイドの埋め込みを一意に決定し、その逆も成り立つことを示している。

Haskellでは、リストデータ構造は自由モノイドだ（いくつかの注意点あり：\urlref{http://comonad.com/reader/2015/free-monoids-in-haskell/}{Dan Doelのブログ記事}を参照）。リスト型\code{{[}a{]}}は、型\code{a}が生成元の集合を表す自由モノイドだ。たとえば、型\code{{[}Char{]}}は単位元――空リスト\code{{[}{]}}――と、\code{{[}'a'{]}}、\code{{[}'b'{]}}などのシングルトン――自由モノイドの生成元――を含む。残りは「積」を適用することで生成される。ここで、二つのリストの積は単純に一方を他方に追加する。追加は結合的かつ単位的だ（つまり、中立元――ここでは空リスト――が存在する）。\code{Char}によって生成された自由モノイドは、\code{Char}の文字からなるすべての文字列の集合に他ならない。Haskellでは\code{String}と呼ばれる：

\src{snippet01}
（\code{type}は型シノニム――既存の型の別名――を定義する。）

別の興味深い例は、一つの生成元から構築された自由モノイドだ。それはユニットのリスト型\code{{[}(){]}}だ。その要素は\code{{[}{]}}、\code{{[}(){]}}、\code{{[}(), (){]}}などだ。そのようなリストはそれぞれ一つの自然数――その長さ――で記述できる。ユニットのリストにはそれ以上の情報は符号化されていない。二つのそのようなリストを追加すると、長さが構成要素の長さの和である新しいリストが得られる。型\code{{[}(){]}}が自然数の加法的モノイド（ゼロあり）と同型であることは容易にわかる。この同型を証明する相互逆関数を示す：

\src{snippet02}
簡単のため、\code{Natural}の代わりに型\code{Int}を使ったが、アイデアは同じだ。関数\code{replicate}は与えられた値――ここではユニット――で事前に満たされた長さ\code{n}のリストを作成する。

\section{いくつかの直感}

以下はいくつかの手を振る議論だ。そのような議論は厳密とは程遠いが、直感を形成するのに役立つ。

自由/忘却随伴について直感を得るには、関手と関数が本質的に情報損失があることを覚えておくと役立つ。関手は複数の対象と射を折りたたむことがあり、関数は集合の複数の要素をまとめることがある。また、それらの像は余域の一部しかカバーしない場合がある。

$\Set$内の「平均的な」hom集合には、最も情報損失が少ないもの（たとえば単射、あるいは同型）から始まり、全域を単一の要素に折りたたむ定数関数（存在する場合）で終わる幅広い関数が含まれる。

任意の圏の射も情報損失があると考える傾向がある。これは精神的なモデルに過ぎないが、随伴について考えるとき――特に片方の圏が$\Set$の場合――には有用なモデルだ。

形式的には、可逆（同型）か非可逆かの射についてしか語れない。後者の種類が情報損失があると考えられる。単射（非折りたたみ）と全射（余域全体をカバー）関数の概念を一般化する単射射と全射射の概念もあるが、単射でも全射でもあり、かつ非可逆な射が存在しうる。

自由 ⊣ 忘却随伴では、左に制約の多い圏$\cat{C}$があり、右に制約の少ない圏$\cat{D}$がある。$\cat{C}$の射は何らかの追加構造を保存しなければならないため「より少ない」。$\cat{Mon}$の場合、乗算と単位を保存しなければならない。$\cat{D}$の射はそれほど多くの構造を保存しなくてよいため「より多く」存在する。

忘却関手$U$を$\cat{C}$の対象$c$に適用するとき、$c$の「内部構造」を明らかにしていると考える。実際、$\cat{D}$が$\Set$の場合、$U$は$c$の内部構造、つまりその台集合を\emph{定義する}と考える。（任意の圏では、対象の内部については他の対象との接続を通じてしか語れないが、ここでは手を振るだけだ。）

$U$を使って二つの対象$c'$と$c$を写すとき、一般に、hom集合$\cat{C}(c', c)$の写像は$\cat{D}(U c', U c)$の部分集合しかカバーしないことが予想される。それは$\cat{C}(c', c)$の射が追加構造を保存しなければならないのに対し、$\cat{D}(U c', U c)$の射はそうしなくてよいからだ。

\begin{figure}[H]
  \centering
  \includegraphics[width=0.45\textwidth]{images/forgettingmorphisms.jpg}
\end{figure}

\noindent
しかし、随伴は特定のhom集合の\newterm{同型}（isomorphism）として定義されるため、$c'$の選択には非常に慎重でなければならない。随伴では、$c'$は$\cat{C}$のどこからでもではなく、（おそらくより小さい）自由関手$F$の像から選ばれる：
\[\cat{C}(F d, c) \cong \cat{D}(d, U c)\]
したがって、$F$の像は任意の$c$への射が豊富な対象で構成されなければならない。実際、$F d$から$c$への構造保存射と$d$から$U c$への非構造保存射が同じ数なければならない。つまり、$F$の像は本質的に構造のない対象（射によって保存すべき構造がない）で構成されなければならない。そのような「構造のない」対象は自由対象と呼ばれる。

\begin{figure}[H]
  \centering
  \includegraphics[width=0.45\textwidth]{images/freeimage.jpg}
\end{figure}

\noindent
モノイドの例では、自由モノイドは単位法則と結合法則によって生成されるもの以外の構造を持たない。それ以外では、すべての乗算は全く新しい要素を生成する。

自由モノイドでは、$2 * 3$は$6$ではなく、新しい要素${[}2, 3{]}$だ。${[}2, 3{]}$と$6$の同一視がないため、この自由モノイドから他のモノイド$m$への射は、それらを別々に写すことができる。しかし、${[}2, 3{]}$と$6$（その積）を$m$の同じ要素に写すことも許される。あるいは加法的モノイドで${[}2, 3{]}$と$5$（それらの和）を同一視することも許される、など。異なる同一視が異なるモノイドを与える。

これは別の興味深い直感につながる：自由モノイドはモノイド演算を実行する代わりに、それに渡された引数を蓄積する。$2$と$3$を乗算する代わりに、$2$と$3$をリストに記憶する。この仕組みの利点は、使用するモノイド演算を指定する必要がないことだ。引数を蓄積し続け、最後にだけ結果に演算子を適用できる。そして、どの演算子を適用するかをその時に選べる。数を加算したり、乗算したり、2を法とする加算を実行したりできる。自由モノイドは式の作成とその評価を分離する。この考えは代数について話すときに再び登場する。

この直感はより複雑な自由構成に一般化される。たとえば、式ツリー全体を評価する前に蓄積できる。このアプローチの利点は、評価をより速くするか、メモリ消費を少なくするためにそのようなツリーを変換できることだ。これは、たとえば行列計算の実装で行われており、熱心な評価は中間結果を格納するための一時配列の多くの割り当てにつながる。

\section{練習問題}

\begin{enumerate}
  \tightlist
  \item
        シングルトン集合を生成元として構築された自由モノイドを考えよ。
        この自由モノイドから任意のモノイド$m$への射と、シングルトン集合から$m$の台集合への関数の間に一対一の対応があることを示せ。
\end{enumerate}
